\documentclass[11pt]{book}

% PRIMARY ENGINE: lualatex essay.tex (two passes mandatory for cross-references)
% FALLBACK ENGINE: xelatex essay.tex (if LuaLaTeX microtype conflict arises)
%
% Known LuaLaTeX issue: microtype may conflict with TeX Gyre Pagella under
% specific font-backend configurations. If LuaLaTeX fails, comment out
% \usepackage{microtype} in preamble.tex and retry, or use XeLaTeX fallback.
% Both engines supported via fontspec; no engine-specific features used.

% preamble.tex
% Shared preamble for "Admissible Reconstruction"
% Engine-agnostic via fontspec (compiles under LuaLaTeX or XeLaTeX)

\usepackage{fontspec}
\setmainfont{TeX Gyre Pagella}

\usepackage{amsmath}
\usepackage{amssymb}
\usepackage{amsthm}
\usepackage{mathtools}
\usepackage[margin=1.25in]{geometry}
\usepackage[colorlinks=true,linkcolor=black,citecolor=black,urlcolor=black]{hyperref}
\usepackage{microtype}
\usepackage{enumitem}
\usepackage{titlesec}

% ----- Theorem environments -----
\theoremstyle{plain}
\newtheorem{theorem}{Theorem}[chapter]
\newtheorem{proposition}[theorem]{Proposition}
\newtheorem{lemma}[theorem]{Lemma}
\newtheorem{corollary}[theorem]{Corollary}

\theoremstyle{definition}
\newtheorem{definition}[theorem]{Definition}
\newtheorem{example}[theorem]{Example}

\theoremstyle{remark}
\newtheorem{remark}[theorem]{Remark}

% ----- Notation macros -----
% Admissibility relation x ~>_R y
\newcommand{\adm}[1]{\rightsquigarrow_{#1}}
% Reachable set of a state
\newcommand{\reach}[1]{R(#1)}
% Time-indexed admissibility structure
\newcommand{\Rt}{R_t}
% Closure functional
\newcommand{\closure}[1]{C(#1)}
% Generativity measure
\newcommand{\genmeas}[1]{G(#1)}
% Observation map
\newcommand{\obsmap}{O}
% History / observable / reconstruction
\newcommand{\History}{H}
\newcommand{\Observable}{O_t}
\newcommand{\Reconstruction}{\hat{H}}

% ----- Section numbering depth -----
\setcounter{secnumdepth}{3}
\setcounter{tocdepth}{2}

% ----- No bulleted lists in body prose: enumitem loaded for the rare
% cases (definitions, notation index) where a list is structurally
% necessary rather than a substitute for prose. -----


\title{\textbf{Structural Invariance Across Modalities}}
\author{Flyxion \\ \textit{Independent Researcher}}
\date{}

\begin{document}

\frontmatter

\maketitle

\chapter*{Abstract}
\addcontentsline{toc}{chapter}{Abstract}

This essay applies the four theorems proven in the monograph
\emph{Admissible Reconstruction} to the projection problem: determining
which structural relations survive when latent intentional states are
realized through different physical channels. Speech, gesture, music, and
audio differ in which aspects of structure each modality preserves and
which it discards, but the formal question remains the same. Does the
projection preserve admissibility structure completely, as required for
representational equivalence, or does it introduce distortions? Does the
projected trajectory maintain nontrivial reachable sets, satisfying the
coherence condition, or does it degenerate? When observers attempt to
reconstruct latent structure from observable projections, does the
observation map produce collisions, and if so, are those collisions
correctable or structural?

The theorems establish that representation, coherence, repair, and
reconstruction are not four independent subjects but four aspects of a
single admissibility relation examined at different points in its
operation. Representation Invariance determines when physically distinct
substrates count as semantically equivalent. Admissibility Preservation
determines when trajectories remain coherent. Monotonic Relaxation
determines when inadmissibility decreases under repair dynamics.
Non-Identifiability determines when reconstruction becomes impossible in
principle, and its diagnostic corollary provides a criterion for
distinguishing correctable observational errors from fundamental
structural limits. Together, these results constrain the space of
possible cross-modal behaviors jointly rather than imposing four
independent constraints.

\tableofcontents

\mainmatter

\chapter{Introduction}
\label{ch:essay-intro}

A spoken word, a gesture, and a musical phrase share no common physical substrate. The acoustic pressure wave that carries speech, the motor trajectory that produces a pointing motion, and the harmonic structure that constitutes a melodic contour differ not merely in details of their realization but in the kind of physical system through which each is instantiated. The puzzle arises when one observes that these heterogeneous signals may nonetheless participate in the communication of related or equivalent meanings. This becomes theoretically puzzling only if the observable signal is assumed to be identical with the structured object being communicated. Once that assumption is rejected, the problem shifts from explaining how different physical substrates could mean the same thing to determining which structural relations survive the projection from latent structure into observable signal, which are altered by the constraints of the physical channel, and which are destroyed so completely that reconstruction becomes impossible.

The thesis developed here is that the observable is a potentially lossy projection of latent structure. The question is therefore not whether reconstruction is possible in general — clearly some projections preserve more structure than others, and some preservation is sufficient for successful communication — but rather what survives projection and what does not. Four theorems developed elsewhere in this research program, proven in the monograph \emph{Admissible Reconstruction}, establish the formal apparatus needed to answer this question precisely. Each theorem addresses a different aspect of the same underlying admissibility relation, examined at different points in its behavior.

The first question is when two representations, possibly realized through entirely different physical or computational substrates, ought to be considered the same object. Representation Invariance, proven in Chapter~1 of \emph{Admissible Reconstruction}, establishes the conditions under which physically different representations count as semantically equivalent. The answer turns on whether the admissibility structures induced by the two representations are related by an appropriate isomorphism, independent of the substrate through which either is realized. The theorem therefore provides a criterion for deciding when substrate differences are semantically irrelevant.

The second question is what it means for a trajectory through admissibility structure to remain coherent. Admissibility Preservation, proven in Chapter~2, addresses this directly. Coherence is not an additional property imposed on top of admissibility; it is what nontriviality of the admissibility relation amounts to when examined along a trajectory rather than at a single state. A trajectory degenerates, and thereby fails to be coherent, exactly when the reachable set collapses in a specific technical sense defined by the theorem. The question is therefore whether the relevant admissibility relation continues to discriminate admissible from inadmissible continuations at each step.

The third question concerns the conditions under which a scalar measure of inadmissibility decreases along a trajectory. Monotonic Relaxation, proven in Chapter~3, establishes that such decrease occurs along trajectories following repair dynamics specifically — transitions selected to increase the volume of the reachable set — but not along admissible trajectories in general. The theorem thereby identifies the precise sense in which certain restorative processes exhibit convergence, while carefully distinguishing this claim from a separate ordering relation developed elsewhere in this program.

The fourth question is when it becomes impossible in principle, not merely difficult in practice, to recover which of two histories actually occurred from the evidence either would have left behind. Non-Identifiability, proven in Chapter~4, establishes both an elementary direction — that when two distinct histories produce identical observables, no single-valued reconstruction function can return the correct distinct history in both cases — and a structural direction, showing that such collisions are not pathological exceptions but a generic feature of observation maps constrained to preserve only reachable-set structure. The theorem additionally provides a diagnostic criterion for distinguishing cases where reconstruction failure reflects a correctable observational error from cases where it reflects a fundamental structural limit.

These four results provide the formal foundation for the cross-modal investigation that follows. Representation, coherence, repair, and the limits of recovery are not four separate phenomena requiring four independent theories. They are four consequences of a single admissibility relation, examined respectively at a single state, along a trajectory, under restorative dynamics, and under lossy projection.

\chapter{Mathematical Foundations}
\label{ch:essay-math}

The four theorems developed in \emph{Admissible Reconstruction} depend on a common apparatus that can be stated concisely. An admissibility structure $R$ on a set of states $X$ determines, for any state $x$, a reachable set $\reach{x}$ consisting of all states $y$ such that $x \adm{R} y$. Nothing about $R$ is assumed beyond this: it is a relation, not a metric, and it need not be symmetric, transitive, or total. The relation may vary with time, written $R_t$, so that what counts as admissible continuation changes as the system evolves. The four theorems below are each derivable from this minimal starting point, with no additional primitive vocabulary required.

\section{Representation Invariance}

Two representations built from different substrates are semantically equivalent when their admissibility structures are related by an admissibility-isomorphism. A map $\varphi : X \to X'$ between state spaces $(X, R)$ and $(X', R')$ is an admissibility-isomorphism when, for all states $x, y \in X$, the biconditional $x \adm{R} y$ if and only if $\varphi(x) \adm{R'} \varphi(y)$ holds. The definition is deliberately minimal: $\varphi$ preserves the admissibility relation in both directions, and nothing more is required.

Theorem~1.1 of \emph{Admissible Reconstruction}, Chapter~1, states that if such a $\varphi$ exists, then $(X, R)$ and $(X', R')$ are semantically equivalent, regardless of the physical or computational substrate realizing either representation. The theorem is close to definitional, and this is by design. Semantic equivalence within the admissibility framework has never been defined as identity of physical encoding but as identity of what a representation permits and forbids as continuation. The theorem makes this definition explicit as an invariance statement rather than leaving it distributed across the apparatus.

What is not definitional is establishing that a particular map $\varphi$ between two representations of interest actually satisfies the biconditional in both directions. This is the harder direction in practice: showing that a transformation preserves reachable-set structure exactly, rather than merely preserving some coarser statistic that happens to look similar. A claim that speech and gesture are representationally invariant, for instance, requires constructing an explicit $\varphi$ and verifying that admissible continuations in one modality correspond exactly to admissible continuations in the other, with neither modality admitting a continuation the other forbids.

Corollary~1.2 of Chapter~1, which is essential to the cross-modal argument developed in later sections, establishes that no property of the state space that is not preserved by every admissibility-isomorphism can be a semantic property of the representation. In particular, substrate-specific properties such as the physical channel through which a representation is realized, its production cost, or its latency are not semantic properties unless they happen to be expressible as constraints on the admissibility relation itself. This does not render substrate properties unimportant; it provides a criterion for deciding when they are semantically relevant. If changing a physical property leaves the admissibility structure invariant, that property does not distinguish the representations at the level recognized by the theory. If changing it alters the reachable relations, then it has become structurally relevant. The corollary thereby draws a sharp boundary between what survives representational equivalence and what belongs to the study of implementation rather than meaning.

\section{Admissibility Preservation}

A trajectory is a sequence of states $S_0, S_1, \dots, S_t$ such that $S_i \adm{R} S_{i+1}$ for each step, where $R$ may itself vary with time. The trajectory degenerates at step $i$ if the reachable set $\reach{S_i}$ is either empty, so that no continuation is admissible, or a singleton entirely determined by the history $S_0, \dots, S_{i-1}$, so that exactly one continuation is admissible and it was already fixed before $S_i$ was reached. In the first degenerate case, the admissibility relation restricted to continuations from $S_i$ is uniformly false; in the second, it is uniformly true of exactly one candidate and uniformly false of every other. Either way, the relation is doing no discriminating work at step $i$.

Theorem~2.1 of Chapter~2 states that a trajectory is coherent if and only if it does not degenerate at any step. Equivalently, the trajectory is coherent if and only if the reachable-set cardinality $|\reach{S_i}|$ exceeds a nontrivial floor relative to what has already been committed at every step along the way. Coherence has sometimes been treated informally as though it were an additional property a trajectory could possess independent of the admissibility relation itself, as though a trajectory could be admissible at every step and yet somehow incoherent, or coherent despite violating admissibility. The theorem rules this out: coherence and nondegeneracy of the admissibility relation are the same property under two names.

Corollary~2.1, proven in the same chapter, addresses operational tests for coherence. Domain-specific measures such as minimum support ratios required of a discourse trajectory, consistency scores, and similar quantities are frequently introduced as though they were primitive definitions of coherence in their own right. The corollary establishes that any such test is correctly understood only as an operational proxy for the condition that $|\reach{S_i}|$ remains above a nontrivial floor, not as an independent criterion. Where a proposed proxy can be false while the reachable set remains nontrivial, or true while the reachable set has degenerated, the proxy has diverged from what coherence actually consists in and should be revised accordingly.

\section{Monotonic Relaxation}

A scalar functional $\closure{X} = -\log \mu(\reach{X})$, where $\mu$ is a measure on the state space assigning nonzero mass to nonempty reachable sets, provides a measure of inadmissibility: the closure functional is large when few futures remain reachable from $X$ and small when many remain. A trajectory follows repair dynamics when each transition $S_i \to S_{i+1}$ is selected, among admissible continuations, to increase reachable-set volume: $\mu(\reach{S_{i+1}}) \ge \mu(\reach{S_i})$. This characterization is inherited from the restorative branch of the modificatory-operation taxonomy developed elsewhere in this research program, where repair is already defined as the operation type that restores lost admissible continuations.

Theorem~3.1 of Chapter~3 states that if a trajectory follows repair dynamics throughout, then the closure functional is nonincreasing along the trajectory: $\closure{S_{i+1}} \le \closure{S_i}$ for every step. The proof is immediate once repair is defined as volume-increasing, since $-\log$ is a strictly decreasing function on positive reals. The theorem's content is therefore not in the derivation but in the definition of repair dynamics being the correct operational characterization. The hypothesis is essential. Constitutive and regulatory modifications, the other two branches of the operation taxonomy, need not increase reachable-set volume, and ordinary state evolution under a fixed admissibility structure can move toward states with strictly smaller reachable sets without this constituting any failure of admissibility. A claim of the form ``inadmissibility always decreases'' would be false. Theorem~3.1 holds only of the restorative subclass, and its scope must be stated with this restriction.

The theorem shares a name with monotonic continuation, an ordering relation $S_t \le S_{t+1}$ developed elsewhere in this program, but the two claims are distinct and should not be conflated. Monotonic continuation is an ordering asserting that later states preserve the distinctions from earlier states that remain necessary for the reachable trajectory; it says nothing about any scalar quantity increasing or decreasing. Monotonic relaxation is a claim that a specific scalar functional decreases along a specific class of trajectories. Corollary~3.1 of Chapter~3 establishes the precise relationship: the two claims coincide only under an additional hypothesis that must be checked independently, and neither implies the other in general. A trajectory can increase reachable-set volume while failing to retain the specific significant continuations required by monotonic continuation, by replacing them with an equal or greater volume of differently located reachable states. Conversely, a trajectory can satisfy monotonic continuation's retention condition while reachable-set volume decreases, if what is retained is a shrinking but still significant core. The corollary exists to prevent the two theorems from being cited interchangeably.

\section{Non-Identifiability}

An observation map $\obsmap : \History \to \Observable$ takes complete histories to whatever observable evidence remains available to an observer without access to the history itself. The elementary direction of the non-identifiability result is nearly definitional. If two distinct histories $H_1 \neq H_2$ satisfy $\obsmap(H_1) = \obsmap(H_2)$, then no function $f : \Observable \to \History$ can satisfy both $f(\obsmap(H_1)) = H_1$ and $f(\obsmap(H_2)) = H_2$ simultaneously, since $f$ must be single-valued. This establishes impossibility once a collision has already been shown to exist, but it says nothing about whether such collisions are rare pathological cases or generic features of realistic observation maps.

The structural direction, Theorem~4.2 of Chapter~4, addresses this directly. Consider any observation map that factors through the reachable-set structure, meaning that $\obsmap(H) = \psi(\reach{H})$ for some function $\psi$, so that the observation map records only which continuations remain admissible from $H$ rather than the full state of $H$ itself. Whenever two histories $H_1 \neq H_2$ satisfy $\reach{H_1} = \reach{H_2}$, it follows immediately that $\obsmap(H_1) = \obsmap(H_2)$, and the elementary direction applies. The theorem's substantive claim is that this is not a degenerate case. The reachable-set operator $R(\cdot)$ partitions the space of histories into equivalence classes by definition: two histories lie in the same class exactly when they share a reachable set. An admissibility structure has a trivial quotient only when no two distinct histories ever share a reachable set, which would require the partition to be discrete. This is not the generic case for admissibility structures where reachable sets are determined by bounded information about the history while the history itself can vary over an unbounded space. The pigeonhole principle then guarantees that histories differing in unrecorded portions while agreeing in recorded portions will share a reachable set and thereby collide under the observation map. Non-injectivity is therefore a structural consequence of the observation map's design, not a contingent misfortune.

Corollary~4.1, the diagnostic corollary, establishes that the question of why reconstruction fails in a given case admits a principled decomposition. Given an observed collision $\obsmap(H_1) = \obsmap(H_2)$ with $H_1 \neq H_2$, one can check whether $\reach{H_1} = \reach{H_2}$. If the reachable sets differ, then the observation map has discarded information that a less lossy procedure could in principle have retained, and the failure is a correctable error of the particular observation procedure. If the reachable sets are equal, then no observation map factoring through reachable-set structure can distinguish $H_1$ from $H_2$, and the failure is a structural limit rather than a procedural defect. The corollary thereby provides a criterion for deciding whether a specific case of reconstruction uncertainty is correctable or fundamental, and this criterion is in principle decidable by examining the reachable-set structure directly.

Together, these four theorems establish that representation, coherence, repair, and the limits of recovery are not four independent subjects requiring separate theoretical machinery. They are four aspects of a single admissibility relation, examined at different points in its operation: Representation Invariance at a single state, Admissibility Preservation along a trajectory, Monotonic Relaxation under restorative dynamics, and Non-Identifiability under lossy projection.

\chapter{Cross-Modal Evidence}
\label{ch:essay-crossmodal}

The formal apparatus developed in the preceding chapter applies whenever structure is projected through a physical channel. Language, gesture, and music differ in which aspects of intentional structure each modality preserves and which it discards, but the projection problem remains the same across all three: determining what survives the transformation from latent structure to observable signal. A worked example illustrates the application of Representation Invariance before turning to measured evidence from the repository's implemented experiments.

Consider the communicative act of requesting transfer of an object, realized in two modalities. In speech, this might be expressed as the utterance ``give me the book,'' consisting of a verb, two pronouns, and a definite noun phrase linearized according to the grammatical constraints of English. In gesture, the same intention might be realized as a sequence consisting of a point toward the speaker, a grasping motion, and a point toward the object, with no single element corresponding uniquely to the verb or to definiteness. The question is whether these two realizations count as representationally invariant in the sense established by Theorem~1.1.

Establishing invariance rigorously requires constructing an explicit admissibility-isomorphism $\varphi$ between the two state spaces and verifying that admissible continuations in one modality correspond exactly to admissible continuations in the other. The speech realization admits continuations such as ``give it to me'' or ``hand me the book,'' which preserve the request structure while varying lexical or syntactic detail. The gesture realization admits continuations such as repeating the sequence with greater emphasis or adding an upward palm orientation to mark urgency. The question is whether there exists a map $\varphi$ such that every admissible speech continuation maps to an admissible gesture continuation and conversely, in both directions, preserving the full reachable-set structure. Without exhibiting such a map and verifying the biconditional, the claim that the two realizations are invariant remains illustrative rather than proven. What can be stated without constructing $\varphi$ is that if such an isomorphism exists, then Corollary~1.2 applies: substrate-specific properties such as the acoustic-versus-motor distinction, production latency, or articulatory cost are not semantic properties unless they constrain the admissibility relation itself.

The question of what information survives projection and what is lost has been investigated directly in the repository's implemented experiments, though not yet across all modalities. The measurements available as of this writing come from experiments on linguistic evolution rather than cross-modal comparison, and they should be understood as evidence concerning projection within a single modality rather than across modalities.

The Phonological Drift experiment, implemented in \texttt{experiments/phonological\_drift.py}, models sound change as a population process rather than as a predetermined historical sequence. Twenty-five speakers occupy a geographical grid, and probabilistic sound changes spread through neighbor influence. Over fifteen generations, the experiment tracks lexical divergence between speakers who began with identical pronunciations. A representative run produces an average lexical divergence of $0.377$ and a maximum divergence of $0.750$ between the most distant speaker pairs. Two speakers pronouncing the ancestral word \textit{pater} may produce \textit{faθer} and \textit{pater} respectively after fifteen generations, despite both trajectories consisting entirely of admissible phonological changes at each step. The measured divergence quantifies how much phonological structure has been altered by the projection through time and population structure, though the experiment does not yet measure how much of the ancestral state remains reconstructible from the divergent outcomes.

The Lexical Natural Selection experiment, implemented in \texttt{experiments/lexical\_selection.py}, models competition among synonymous forms without reducing fitness to a universal scalar. Variants differ along dimensions including length, articulatory cost, regularity, and prestige, with relative advantages depending on context. Four meanings begin with three competing variants each and evolve over one hundred fifty time steps. A representative run for the meaning \textit{father} produces a distribution in which the variant \textit{pa} accounts for $36.5\%$ of usage, \textit{papa} for $34.8\%$, and \textit{pater} for $28.7\%$. The experiment demonstrates that selection can alter the frequency distribution of competing forms without eliminating variation entirely, though again it does not yet measure how much of the selection history remains recoverable from the final distribution alone.

The Reconstruct experiment, implemented in \texttt{experiments/reconstruct.py}, closes the experimental loop by attempting historical inference without access to the simulator's privileged ground truth. A proto-language is evolved into four daughter languages over twenty time steps, the ancestral state is concealed, and a simplified comparative method attempts reconstruction using only the surviving daughters. A representative run achieves correct reconstruction for six of ten test items, partial reconstruction for four items, and complete failure for zero items, yielding a sixty-percent exact-match accuracy. This is the first direct measurement of reconstruction accuracy under the complete $\History \to \Observable \to \Reconstruction$ protocol, and it demonstrates that some information survives the projection from history to observable well enough to permit reconstruction while other information is lost.

The Recoverability experiment, implemented in \texttt{experiments/recoverability.py}, investigates the distinction between reconstruction error and fundamental non-identifiability. Rather than asking only whether a particular history can be reconstructed, the experiment generates multiple independent histories and searches for cases in which $H_1 \neq H_2$ but $\obsmap(H_1) \approx \obsmap(H_2)$. A representative run across ten generated histories identifies twenty-four pairs of distinct histories with similar observable outcomes, including cases where observable distance is measured as zero despite the histories differing in event count and evolutionary trajectory. These cases instantiate the structural direction of Theorem~4.2: distinct histories sharing the same reachable-set structure produce indistinguishable observables, and no reconstruction procedure operating solely on observables can recover which history actually occurred.

The measurements reported above are all drawn from experiments operating within a single modality, specifically linguistic evolution under sound change, lexical competition, and historical divergence. Cross-modal measurements comparing information loss across language, gesture, music, and audio are not yet available from implemented experiments. Claims about the percentage of trajectory information lost when continuous gesture is sampled to discrete frames, or the percentage of semantic content lost when structured meaning is linearized into speech, remain hypotheses to be tested rather than measured results. A proper cross-modal comparison would require implementing projectors and inverse engines for each modality as specified in the unified framework, running controlled experiments in which the same semantic structure is projected through different channels, and measuring what fraction of the original structure remains recoverable in each case. Such experiments would test whether the information-loss patterns predicted by the four theorems manifest differently across modalities or whether the projection problem exhibits modality-independent structure.

The evidence available from the repository's current experiments demonstrates that the $\History \to \Observable \to \Reconstruction$ framework can distinguish three questions that are easily conflated: what evolutionary history actually occurred, what evidence of that history survived to the present, and what an observer can legitimately infer from that surviving evidence. The Reconstruct experiment shows that reconstruction can succeed partially even when information has been lost. The Recoverability experiment shows that distinct histories can become observationally indistinguishable, instantiating the collision condition required by Theorem~4.2. Together, these results establish that the framework's epistemic distinctions are not merely definitional but measurable, and that the theorems proven in \emph{Admissible Reconstruction} apply to the implemented experiments. What remains to be measured is whether the same distinctions apply with comparable quantitative characteristics when the projection operates across modalities rather than solely through time within a single modality.

\chapter{Perceptual Control Theory}
\label{ch:essay-pct}

The four theorems established in Chapter~\ref{ch:essay-math} describe structural properties of the admissibility relation without specifying any particular mechanism for how behavioral output $B_t$ is generated from control states or intentional structures. Perceptual Control Theory, as developed by Powers and others, offers one candidate mechanism for that generation process, though it is important to state at the outset that PCT is compatible with the admissibility framework rather than derived from it. The theorems concerning representation, coherence, repair, and non-identifiability remain valid independent of whether any particular cognitive or behavioral system operates according to control principles, and the empirical hypotheses developed in this chapter require independent validation rather than following as mathematical consequences of the formal apparatus.

The central claim of Perceptual Control Theory, as articulated by Powers, is that behavior should be understood not as a response emitted in reaction to external stimuli but as the continuously variable output of a negative-feedback process attempting to maintain selected perceptions near internally specified reference states. The fundamental control loop can be written schematically as a sequence in which a reference value $r$ is compared with a currently perceived value $p$ to produce an error signal $e = r - p$, which drives behavioral output $o$, which acts on the environment $E$, which in turn affects subsequent perception, closing the loop. What an organism controls, in this framework, is not its behavioral output directly but rather the perceptual consequences of that output. The behavior itself is therefore free to vary across a potentially wide range whenever different motor actions succeed equally well in maintaining the controlled perception near its reference value.

If linguistic behavior is generated according to control principles rather than as a fixed stimulus-response mapping, several empirical predictions follow, though each requires stating as a hypothesis rather than as an established result. First, if speakers control perceptions such as listener comprehension, social alignment, or production effort rather than controlling linguistic forms directly, then the same communicative intention should be realizable through multiple distinct surface forms whenever those forms maintain the relevant controlled variables equally well. This predicts that synonymy, paraphrase, and stylistic variation should be widespread in natural language, which is consistent with observation but does not by itself establish that control is the generating mechanism. Second, if the reference values for controlled perceptions can be adjusted dynamically in response to environmental or social conditions, then linguistic behavior should adapt to context in ways that a fixed grammar or lexicon would not predict. For instance, if ambient noise in the environment reduces the probability that phonetic distinctions are perceived correctly, a control-based account predicts that speakers should adjust their articulation to maintain comprehensibility despite the degraded channel, which could manifest as increased articulatory precision, greater reliance on redundant cues, or shifts toward more acoustically distinct phonemes. Third, if different controlled perceptions can demand conflicting adjustments, then linguistic behavior should exhibit trade-offs rather than optimizing any single dimension. A speaker might simultaneously experience pressure to increase explicitness for comprehension, reduce articulatory effort for ease of production, and conform to community norms for social alignment, with no single adjustment satisfying all three demands optimally. The resolution of such conflicts could occur through fixed hierarchical priority, dynamic adjustment of reference levels, or independent controllers operating at different linguistic levels, and each resolution mechanism predicts a different pattern of observable linguistic behavior.

The connection between control-based generation of behavior and the formal apparatus of admissibility operates through Theorem~2.1, Admissibility Preservation. Whatever mechanism generates the behavioral trajectory $B_t$, whether control-based or otherwise, the resulting trajectory through state space remains subject to the coherence condition established by that theorem. A trajectory degenerates, and thereby fails to be coherent, when the reachable set collapses to the empty set or to a singleton already determined by prior history, regardless of how the transitions were selected. A control process that drives the system into a state from which no admissible continuation exists has produced an incoherent trajectory in the technical sense defined by the theorem, even if the control loop itself remains operationally stable. The admissibility framework therefore constrains which trajectories count as coherent independent of the generative mechanism, while control theory offers one account of how a system might navigate through the space of admissible trajectories without prespecifying which particular path will be taken.

The most direct connection between PCT and the projection problem addressed in this essay concerns the equivalence class of behaviors that maintain a given controlled perception. If multiple distinct surface forms can maintain listener comprehension at the same level, those forms occupy an equivalence class with respect to the comprehension controller, even if they differ substantially in other dimensions such as phonetic detail, syntactic structure, or lexical choice. This is precisely the structure addressed by Representation Invariance: if the admissibility structure induced by the comprehension constraint is preserved across the different surface forms, then those forms count as semantically equivalent in the sense established by Theorem~1.1, despite their substrate differences. The hypothesis that natural linguistic systems operate according to control principles therefore predicts that representational equivalence classes should align with the equivalence classes defined by controlled perceptual variables, which is empirically testable by measuring whether forms judged semantically equivalent by speakers maintain the same controlled variables and whether forms that differ in controlled variables are judged semantically distinct.

The ecological examples mentioned earlier illustrate how control-based adaptation might interact with environmental constraints to produce linguistic change, though again each example requires marking as a hypothesis rather than an established mechanism. If ambient noise degrades the channel through which speech is transmitted, and if speakers control for listener comprehension rather than for preservation of particular phonetic forms, then the prediction is that sound systems should adapt toward greater robustness under the specific noise conditions encountered. In a maritime environment with wind and wave noise concentrated at low frequencies, this might favor consonant systems with more high-frequency distinctions. In an urban environment with broad-spectrum mechanical noise, it might favor vowel systems with greater formant separation or increased reliance on tonal contrasts less susceptible to spectral masking. The same environmental condition encountered by a population controlling different perceptual variables, such as articulatory ease rather than comprehensibility, would predict a different adaptive trajectory. Critically, these predictions concern the direction and character of adaptation rather than merely asserting that adaptation occurs, and they are in principle testable by comparing sound-system evolution across populations in documented acoustic environments, though such data are not yet available in the repository's implemented experiments.

The relationship between control theory and the Non-Identifiability theorem merits explicit statement. If two populations arrive at similar observable linguistic states through different control histories, one driven primarily by environmental noise favoring robust contrasts and another driven primarily by social differentiation producing superficially similar forms, then Theorem~4.2 applies: an observer examining only the contemporary languages without access to the complete control history faces a potential collision in which $\obsmap(H_1) = \obsmap(H_2)$ despite $H_1 \neq H_2$. The diagnostic corollary, Corollary~4.1, provides a criterion for assessing whether such a case represents a structural limit or a correctable observational error by checking whether the two histories share the same reachable-set structure. If the reachable sets differ, additional evidence could in principle distinguish the histories. If they are equal, no observation map factoring through reachable-set structure can resolve the ambiguity. The hypothesis that linguistic systems are generated by control processes thereby introduces a specific class of potential collisions in which different control conditions produce the same linguistic outcomes, and the theorem establishes which such collisions are fundamentally irresolvable.

These connections between control theory and the admissibility framework remain hypotheses requiring empirical validation. The theorems proven in \emph{Admissible Reconstruction} establish what follows from the admissibility relation itself, independent of any commitment to control as the generative mechanism. Control theory offers a particular account of how behavioral trajectories might be generated in a way that interacts naturally with admissibility constraints, equivalence classes, and observational limits, but that interaction is a conjecture about empirical systems rather than a mathematical consequence of the formal apparatus. The value of stating the connection explicitly is that it generates testable predictions concerning equivalence classes, environmental adaptation, conflict resolution, and historical non-identifiability that can be investigated through targeted experiments rather than remaining at the level of informal speculation.

\chapter{Theoretical Synthesis}
\label{ch:essay-synthesis}

Representation Invariance, Admissibility Preservation, Monotonic Relaxation, and Non-Identifiability are not four independent results requiring separate theoretical foundations. They are four consequences of a single admissibility relation examined at different points in its operation: Theorem~1.1 at a single state, Theorem~2.1 along a trajectory, Theorem~3.1 under restorative dynamics, and Theorem~4.2 under lossy projection. What varies is not the fundamental object of study but the operation being performed on it. This unity becomes consequential when one recognizes that the same structural mechanism can manifest as apparently different phenomena depending on where in the system it is observed.

Consider the reachable-set collapse that defines degeneracy in Theorem~2.1. A trajectory becomes incoherent when $\reach{S_i}$ collapses to the empty set or to a singleton already determined by prior history, because at that point the admissibility relation ceases to discriminate among continuations in any non-vacuous way. This is a forward-looking failure: the system has reached a state from which it cannot continue coherently. The same collapse, when examined retrospectively rather than prospectively, produces the observational collisions addressed by Theorem~4.2. If two distinct historical trajectories $H_1$ and $H_2$ both lead to states sharing the same reachable-set structure, then any observation map that records only what remains admissible from those states cannot distinguish which trajectory actually occurred. The forward-looking degeneracy and the backward-looking non-identifiability are therefore not separate failures but the same reachable-set structure examined from opposite temporal directions.

A concrete illustration clarifies this relationship. Suppose a discourse trajectory in which a speaker's utterances progressively narrow the space of admissible next contributions until only a single continuation remains licensed by what has been said so far. At that point the trajectory has degenerated in the sense of Theorem~2.1, because $\reach{S_t}$ is a singleton already fixed by $S_0, \dots, S_{t-1}$, and the admissibility relation is no longer performing discriminative work. Now consider an observer attempting to reconstruct how the discourse arrived at state $S_t$ from evidence consisting solely of which continuations remain admissible at that point. If two different discourse histories, one that narrowed possibilities through explicit assertion and another that narrowed them through presupposition and implicature, both produce the same singleton reachable set, then the observation map $\obsmap(H) = \reach{S_t}$ produces a collision: $\obsmap(H_1) = \obsmap(H_2)$ despite $H_1 \neq H_2$. The diagnostic corollary, Corollary~4.1, then applies: since the two histories share the same reachable set by construction, no observation map factoring through reachable-set structure can distinguish them, and the reconstruction failure is a structural limit rather than a correctable observational error.

The joint treatment thereby reveals that what appears as incoherence when examined forward through a trajectory and what appears as non-identifiability when examined backward through a reconstruction are related manifestations of reachable-set structure. The distinction is not merely that the same mathematical apparatus applies to both cases; it is that the collapse of discriminative structure causes both phenomena simultaneously. A system that has degenerated into a state from which only one continuation is admissible has also, by that very fact, entered a state from which its past becomes harder to recover, because the richness of the reachable set is precisely what would have allowed an observer to infer how that state was reached. Conversely, a pair of histories that produce observationally indistinguishable outcomes do so because they lead to states with identical reachable structure, and identical reachable structure is exactly the condition under which the forward trajectory from either state would be indistinguishable.

This relationship extends to the repair dynamics addressed by Theorem~3.1. Repair, defined as transitions that increase reachable-set volume, moves the system away from degeneracy by expanding $\mu(\reach{S_t})$. Since degeneracy is the condition under which both coherence fails and historical non-identifiability becomes more likely, repair trajectories that successfully increase reachable-set volume thereby simultaneously restore coherence in the forward direction and create observational distinctions that make historical reconstruction more feasible in the backward direction. The scalar energy functional $\closure{S_t} = -\log \mu(\reach{S_t})$ therefore measures not only inadmissibility in an abstract sense but also, concretely, the degree to which the system has approached the dual failures of forward incoherence and backward non-identifiability. Monotonic decrease of $\closure{S_t}$ along a repair trajectory signals movement away from both failure modes at once.

The cross-modal application developed in this essay benefits from this joint treatment because projection across modalities introduces opportunities for reachable-set structure to be preserved, altered, or destroyed in ways that affect representation, coherence, and recoverability simultaneously. If projecting semantic structure into speech preserves certain reachable-set relations while projecting the same structure into gesture preserves different relations, then Theorem~1.1 establishes that the two projections are representationally equivalent only to the extent that the preserved relations coincide. If either projection produces states with collapsed reachable sets, then Theorem~2.1 predicts that continuations in that modality will degenerate. If both projections lead to states sharing the same reachable structure despite originating from different semantic intentions, then Theorem~4.2 predicts that reconstructing the original intention from the projected signal will encounter non-identifiability. The four theorems thereby constrain the space of possible cross-modal behaviors jointly rather than imposing four independent constraints that must be checked separately.

Treating the four theorems as aspects of a single admissibility relation does not reduce them to triviality; rather, it allows their interactions to be studied systematically. A mechanism that generates behavior, whether control-based as discussed in Chapter~\ref{ch:essay-pct} or otherwise, operates within a space of trajectories that must satisfy the coherence condition of Theorem~2.1 if they are to avoid degeneracy. The observable projection of those trajectories is subject to the non-identifiability condition of Theorem~4.2 whenever the projection discards distinctions beyond those encoded in reachable-set structure. The representational equivalences licensed by Theorem~1.1 determine which different physical realizations of behavior count as semantically interchangeable. And the repair dynamics of Theorem~3.1 provide one account of how systems might navigate toward regions of state space where reachable sets remain nontrivial, thereby avoiding the dual failures of incoherence and non-identifiability. These are not four separate constraints applied in sequence but four views of the same structural condition examined from different angles, and recognizing them as such clarifies both what the admissibility framework establishes and what it leaves open for empirical investigation.

\chapter{Implications}
\label{ch:essay-implications}

The theorems proven in \emph{Admissible Reconstruction} establish consequences that follow necessarily from the admissibility framework and generate hypotheses that require independent empirical validation. The distinction between these two categories must be maintained throughout.

\section{Established Consequences}

Theorem~4.2, the structural direction of Non-Identifiability proven in Chapter~4, establishes that observation maps constrained to preserve only reachable-set structure are generically non-injective whenever the space of histories exceeds what that structure can encode. This applies directly to historical linguistics. Sound change operating within a language lineage and lexical borrowing from a neighboring language can produce phonologically identical outcomes in the descendant language despite originating from entirely different historical processes. If two evolutionary histories $H_1$ and $H_2$, one involving regular sound change and the other involving contact-induced borrowing, both lead to the same phonological form in the observable contemporary language, and if that phonological form determines the same set of admissible continuations regardless of which history produced it, then the two histories share the same reachable-set structure: $\reach{H_1} = \reach{H_2}$. Under these conditions, Theorem~4.2 guarantees that any observation map $\obsmap$ factoring through reachable-set structure produces a collision: $\obsmap(H_1) = \obsmap(H_2)$. No reconstruction procedure operating solely on the observable contemporary form can determine with certainty which history actually occurred.

The diagnostic corollary, Corollary~4.1, provides a criterion for assessing such cases. An observer examining a contemporary form and suspecting that it could have arisen either through inheritance with regular sound change or through borrowing from a contact language can, in principle, check whether the two candidate histories would produce the same reachable set. If the reachable sets differ, then additional evidence—perhaps from comparative data, geographical distribution, or semantic fields—could distinguish the histories, and the reconstruction uncertainty is correctable. If the reachable sets are equal, then the uncertainty is a structural limit: no observation map preserving only what remains admissible from the contemporary state can resolve the ambiguity, regardless of how much comparative data is gathered or how sophisticated the reconstruction algorithm becomes. This is not a claim about the practical difficulty of reconstruction but about what information has ceased to exist in the observable projection.

The same consequence applies to any domain in which distinct processes can produce observationally equivalent outcomes. Two musical traditions that arrive at similar harmonic structures, one through independent evolution within a culture's aesthetic constraints and another through cultural contact and adoption, face the same non-identifiability condition if the contemporary harmonic structure determines the same space of admissible continuations regardless of origin. Two gestural systems that develop similar pointing conventions, one through communicative necessity in a specific ecological niche and another through deliberate pedagogical standardization, cannot be distinguished by an observer who has access only to the contemporary gestural repertoire if both produce the same reachable-set structure. In each case, Theorem~4.2 establishes that the observational collision is a property of the projection map rather than a failure of the observer's inferential capacity.

\section{Empirical Hypotheses}

The framework developed in this essay generates several testable predictions concerning cross-modal cognition, though each requires marking explicitly as a hypothesis rather than an established result.

If gesture recognition in human cognition operates as an inverse reconstruction problem attempting to recover latent embodied semantic structure from observed motor trajectories, as discussed in Chapter~\ref{ch:essay-intro}, then the hypothesis is that recognition systems incorporating temporal trajectory information should outperform systems operating solely on isolated frames. The formal motivation is that continuous motor trajectories preserve reachable-set structure that is lost when the trajectory is sampled to discrete configurations. A recognition system that attempts to reconstruct the latent intention by inverting the projection from intention through motor planning to observable configuration should therefore perform better when given access to the continuous trajectory than when given only a sequence of static poses, because the trajectory preserves distinctions that the frame sequence discards. This is testable by comparing recognition accuracy for trajectory-aware versus frame-based algorithms on the same gesture corpus, though such data are not available in the repository's current experiments.

If musical cognition involves reconstructing structured temporal relationships from acoustic projections, then a second hypothesis is that listeners should exhibit systematic reconstruction errors in cases where distinct compositional structures project to similar acoustic surfaces. Specifically, if two different harmonic progressions produce similar spectral and temporal patterns but differ in their reachable continuations—what harmonic moves remain admissible next—then listeners should be more accurate at predicting continuations when given the full compositional structure than when given only the acoustic surface. This predicts measurable differences in continuation judgment tasks depending on whether the task provides structural or surface information, and it further predicts that the error patterns should align with cases where Theorem~4.2 applies: listeners should exhibit greater uncertainty exactly when the observation map produces collisions between distinct compositional histories.

If control processes in biological systems maintain reference values for perceptual variables rather than for output forms, as discussed in Chapter~\ref{ch:essay-pct}, then a third hypothesis is that neural correlates of linguistic behavior should exhibit sustained activity patterns associated with controlled perceptual goals rather than transient patterns associated with individual motor outputs. This predicts that neuroimaging studies should find persistent activation in regions associated with comprehension monitoring, social feedback processing, and effort evaluation throughout utterance production, rather than activation patterns that rise and fall with individual articulatory gestures. The prediction is testable through targeted neuroimaging protocols, though it requires careful experimental design to distinguish sustained goal-related activity from preparatory or anticipatory motor activity, and it remains an empirical hypothesis rather than a consequence of the formal framework.

The distinction between the established consequences in the preceding section and these empirical hypotheses is fundamental. Theorem~4.2 proves that observation maps factoring through reachable-set structure are non-injective under the stated conditions, and this applies to any domain satisfying those conditions, including historical linguistics. The theorem does not prove that human gesture recognition operates as an inverse reconstruction problem, that musical cognition involves reachable-set recovery, or that linguistic behavior is generated by perceptual control loops. Each of these is an empirical conjecture motivated by the formal framework but requiring independent validation through targeted experiments, behavioral measurements, or neurophysiological data. The value of stating them explicitly as hypotheses is that they become testable predictions with clear success and failure conditions rather than remaining at the level of informal speculation, and they illustrate how the admissibility framework can generate empirically consequential predictions across modalities while respecting the boundary between what the mathematics establishes and what the world must be observed to determine.

\chapter{Fundamental Limits}
\label{ch:essay-limits}

The Non-Identifiability theorem establishes that certain failures of historical reconstruction are not contingent limitations of particular inference methods but structural properties of the observation map itself. This chapter develops that result in detail, with particular attention to the equivalence-class structure that makes non-identifiability inevitable under specified conditions and to the diagnostic criterion that distinguishes structural limits from correctable observational errors.

An observation map $\obsmap : \History \to \Observable$ transforms complete histories into whatever evidence remains accessible to an observer at a given time. The map is deliberately lossy: it discards information that was present in the history but has not survived to the observable. The question is which histories can be distinguished on the basis of observables alone and which cannot. Two histories $H_1$ and $H_2$ are observationally indistinguishable when $\obsmap(H_1) = \obsmap(H_2)$ despite $H_1 \neq H_2$. Such histories form an equivalence class under the observation map, written $[H]_{\obsmap} = \{ H' : \obsmap(H') = \obsmap(H) \}$, consisting of all histories that produce the same observable outcome as $H$. The size and structure of these equivalence classes determine what can and cannot be recovered through reconstruction.

The elementary direction of the non-identifiability result, Theorem~4.1 of Chapter~4 in \emph{Admissible Reconstruction}, is immediate. If $H_1$ and $H_2$ are distinct members of the same equivalence class, so that $\obsmap(H_1) = \obsmap(H_2)$, then no function $f : \Observable \to \History$ can return both $H_1$ and $H_2$ as the correct reconstruction of the single observable $\obsmap(H_1) = \obsmap(H_2)$, because functions must be single-valued. Any reconstruction procedure faced with this observable must choose one history or the other, or report that the evidence is insufficient to decide, but it cannot return both as the unique correct answer. This establishes impossibility once a collision has been demonstrated, but it does not yet explain why such collisions should be expected to occur or whether they are rare pathologies confined to contrived examples.

The structural direction, Theorem~4.2, addresses this directly and constitutes the deepest result concerning reconstruction limits. Consider observation maps that factor through the reachable-set structure, meaning that the observable is determined by which continuations remain admissible from the history rather than by the full details of the history itself. Formally, $\obsmap(H) = \psi(\reach{H})$ for some function $\psi$, where $\reach{H}$ is the reachable set from the terminal state of history $H$. Such maps record what futures are possible from the current state but not which particular path was taken to arrive there. This class of observation maps is not artificially restrictive; it includes any observation procedure that determines the observable solely by examining what remains admissible continuation rather than by recovering the complete temporal trace of how the system evolved.

Under this condition, whenever two histories $H_1 \neq H_2$ satisfy $\reach{H_1} = \reach{H_2}$, it follows immediately that $\obsmap(H_1) = \psi(\reach{H_1}) = \psi(\reach{H_2}) = \obsmap(H_2)$, producing a collision. The substantive claim is that such collisions are not exceptional. The reachable-set operator $R(\cdot)$ partitions the space of histories into equivalence classes by definition: histories lie in the same class exactly when they share the same reachable set. An admissibility structure induces a trivial partition, in which every equivalence class is a singleton, only when no two distinct histories ever produce the same reachable set. This would require that the reachable set encode the full history rather than merely the current admissible-continuation structure, which contradicts the premise that the observation map factors through reachable sets precisely because it discards historical detail beyond what remains relevant for future admissibility.

The generic case is therefore non-trivial equivalence classes. When histories are individuated by more information than their reachable sets preserve, the pigeonhole principle applies: there are more histories than distinct reachable-set structures, so multiple histories must map to the same structure. These are not edge cases or carefully constructed counterexamples but the expected outcome of any observation map that preserves continuation structure without preserving the full temporal trace of how that structure was established. The structural direction thereby establishes that non-injectivity is a property of the observation map's design rather than a correctable defect.

The equivalence-class perspective clarifies what reconstruction can and cannot achieve. A reconstruction algorithm operating on observables $o = \obsmap(H)$ without access to the hidden history $H$ can, at best, identify the equivalence class $[H]_{\obsmap}$ to which the true history belongs. If that class contains only one element, reconstruction can succeed uniquely. If it contains multiple elements, reconstruction faces an inherent ambiguity: every member of the class is consistent with the observable, and no additional evidence drawn from the same observation map can distinguish them. The question is then whether the ambiguity can be resolved by switching to a less lossy observation map or whether it persists even when the observation preserves all available continuation structure.

This is where the diagnostic corollary, Corollary~4.1 of Chapter~4, becomes essential. Given a collision $\obsmap(H_1) = \obsmap(H_2)$ with $H_1 \neq H_2$, the corollary establishes a criterion for diagnosing the source of reconstruction failure. The diagnosis proceeds by checking whether $\reach{H_1} = \reach{H_2}$. If the reachable sets differ, then the two histories induce different continuation structures, and the observation map has discarded information that a less lossy procedure could in principle have retained. The failure is correctable: a different observation map that preserves more of the reachable-set structure, or that records additional features beyond reachable sets, could distinguish the histories. The reconstruction error is therefore a limitation of the particular observation procedure employed rather than a fundamental property of the histories themselves.

If, conversely, $\reach{H_1} = \reach{H_2}$, then the two histories produce identical continuation structures, and no observation map that factors through reachable-set structure can distinguish them. The equivalence class $[H_1]_{\obsmap}$ contains at least $H_1$ and $H_2$ as distinct elements, and this is not an artifact of the particular $\psi$ chosen but a structural consequence of the histories sharing the same reachable set. Under these conditions, the reconstruction failure is a fundamental limit: the information required to distinguish the histories has ceased to exist in the observable continuation structure, and no refinement of the observation procedure that preserves the factorization through reachable sets can recover it.

The diagnostic corollary thereby provides a principled answer to the question of whether a specific case of reconstruction uncertainty reflects insufficient evidence, inadequate inference methods, or genuine non-recoverability. The distinction is in principle decidable by examining the reachable-set structure directly. In practice, this requires access to the reachable sets $\reach{H_1}$ and $\reach{H_2}$, which may not be directly observable, but the theoretical criterion is clear: check whether the continuation structures coincide. If they do, the limit is structural. If they do not, additional evidence could resolve the ambiguity.

The theorem has broader implications beyond its immediate statement. It establishes that the equivalence classes $[H]_{\obsmap}$ are not arbitrary collections of histories that happen to produce similar outcomes by coincidence but are precisely the sets of histories sharing the same reachable-set structure. The size and internal structure of these classes are therefore determined by the relationship between the space of possible histories and the space of reachable-set structures that those histories can induce. When history space is high-dimensional and reachable-set space is lower-dimensional, equivalence classes are necessarily non-trivial, containing multiple distinct histories. The observation map then performs a many-to-one projection, and reconstruction from observables becomes inherently ambiguous for any history whose equivalence class contains more than one element.

This is the deepest sense in which observables are lossy projections of latent structure. The loss is not merely that some details are omitted, as though a sufficiently careful observer could in principle recover them by examining the evidence more closely. The loss is that entire dimensions of variation in the history space are collapsed into equivalence classes in the observable space, and no observation map preserving only reachable-set structure can recover those collapsed dimensions. The structural non-identifiability established by Theorem~4.2 is therefore not a limitation of observers but a property of the channel through which history is projected into observables. Different histories that induce the same continuation structure have become, from the perspective of any observation map factoring through that structure, the same observable outcome, and the distinction between them is no longer present in the evidence to be discovered.

\chapter{Objections}
\label{ch:essay-objections}

A reasonable objection to the framework developed in this essay is that historical linguistics has long acknowledged that some reconstructions remain uncertain and that not all aspects of linguistic history can be recovered with equal confidence. Comparative linguists distinguish between secure reconstructions supported by regular correspondences across multiple daughter languages and speculative reconstructions based on limited or ambiguous evidence. The field has developed explicit notation for marking reconstruction confidence, with asterisks indicating unattested reconstructed forms and question marks or parentheses indicating uncertain elements. Methodological discussions routinely address the limitations of the comparative method, acknowledging that sound change can obscure original distinctions, that borrowing can introduce forms that resist genealogical explanation, and that chance resemblances can produce spurious cognate sets. Given this existing recognition of reconstruction uncertainty, the objection runs, what does formalization through the admissibility framework add beyond restating in mathematical notation what practitioners already know from experience?

The objection gains force from observing that the practical work of historical reconstruction has proceeded successfully for two centuries without requiring theorems about non-identifiability or diagnostic criteria for distinguishing structural limits from correctable errors. Indo-European comparative linguistics established sound laws, reconstructed proto-forms, and built phylogenetic trees before the formal apparatus developed in this essay existed. The comparative method identifies cognates, projects ancestral states, and assesses reconstruction confidence through procedures that are well understood and widely applied without reference to reachable-set structures or equivalence classes under observation maps. If the field's existing methodological toolkit already handles reconstruction uncertainty effectively, then formalizing that uncertainty as a theorem about observation maps factoring through reachable-set structure might appear to be an exercise in restating the obvious in unnecessarily technical language. The objection deserves serious consideration because it correctly identifies that the general fact of reconstruction uncertainty is not a novel discovery.

The answer to this objection must begin by acknowledging what is correct in it. The theorem does not establish that historical reconstruction is sometimes uncertain; that fact is indeed already known and has been incorporated into linguistic methodology for generations. Theorem~4.2 does not tell historical linguists anything they did not already understand qualitatively about the difficulty of distinguishing sound change from borrowing, the challenges posed by sporadic changes, or the problems created by lexical replacement. The formalization is not intended as a corrective to a field that was previously unaware of these issues. What the theorem provides, and what existing methodology does not, is a principled criterion for diagnosing why a particular reconstruction is uncertain and whether that uncertainty is correctable or fundamental.

The contribution lies specifically in Corollary~4.1, the diagnostic corollary developed in Chapter~4 of \emph{Admissible Reconstruction} and discussed in detail in Chapter~\ref{ch:essay-limits}. The corollary establishes that given a collision between two candidate histories $H_1$ and $H_2$ producing the same or similar observables, one can in principle check whether $\reach{H_1} = \reach{H_2}$. If the reachable sets differ, the reconstruction uncertainty reflects information that the particular observation procedure employed has discarded but that a less lossy procedure could retain. The failure is correctable: gathering different kinds of evidence, applying different comparative techniques, or examining features beyond those currently observed could distinguish the histories. If the reachable sets are equal, the reconstruction uncertainty reflects information that has ceased to exist in any observation preserving only continuation structure. The failure is a structural limit: no amount of additional evidence of the same kind, and no refinement of comparative methodology operating on the same observable features, can resolve the ambiguity.

This distinction is not drawn by existing methodology. When a historical linguist encounters an uncertain reconstruction, current practice might note that the evidence is ambiguous, that multiple hypotheses are consistent with the data, or that the reconstruction should be marked as tentative. What current practice does not provide is a criterion for determining whether the ambiguity is in principle resolvable through additional evidence or whether it reflects a genuine limit of what the observable projection preserves. The difference is not merely academic. If an uncertain reconstruction reflects a correctable observational gap, then the appropriate response is to seek additional evidence: examine more daughter languages, investigate dialectal variation, or incorporate non-phonological data such as morphological alternations or semantic distributions. If it reflects a structural limit, then such efforts are unlikely to succeed, and the appropriate response is to acknowledge the ambiguity as irresolvable rather than treating it as a temporary gap awaiting future discovery.

Consider a concrete example. Suppose a historical linguist reconstructing Proto-Polynesian encounters a lexical item that appears in some daughter languages with initial \textit{p} and in others with initial \textit{f}, and suppose two hypotheses are proposed: either the proto-form had \textit{p} with subsequent lenition to \textit{f} in some branches, or the proto-form had \textit{f} with subsequent fortition to \textit{p} in others. Current methodology would examine the broader pattern of sound correspondences, look for conditioning environments, and assess which hypothesis better fits the overall phonological system, but it might conclude that the evidence is insufficient to decide with certainty. The diagnostic corollary provides an additional step: check whether the two hypothesized histories lead to states with the same reachable-set structure. If both reconstructed proto-forms license the same set of admissible phonological developments in the daughter languages, then the two histories share a reachable set, and Theorem~4.2 applies: no observation map factoring through phonological continuation structure can distinguish them. The uncertainty is structural. If the reachable sets differ—if, for instance, reconstructing \textit{p} predicts different continuation patterns in unstudied daughter languages than reconstructing \textit{f}—then the uncertainty is correctable, and the appropriate response is to examine those unstudied languages or other evidence that would fall within the differential reachable sets.

The objection is therefore correct that the field already knows reconstruction is sometimes uncertain, but it underestimates what the formalization provides. The theorem does not merely restate that uncertainty exists; it provides a diagnostic tool for classifying uncertainty into correctable and structural categories, and it establishes that this classification is in principle decidable by examining reachable-set structure. This transforms reconstruction uncertainty from an undifferentiated acknowledgment that evidence may be insufficient into a structured question with a principled answer: is this case uncertain because we have not yet gathered the right evidence, or because the evidence required to decide has ceased to exist in the observable projection? The formalization thereby adds precision to a distinction that existing methodology gestures toward but does not make explicit or decidable.

\chapter{Future Directions}
\label{ch:essay-future}

Future work divides into extensions that strengthen the applicability of existing theorems by providing additional empirical measurements and extensions that require new theoretical machinery beyond the admissibility framework developed here. The first category includes implementing projector and inverse-engine pairs for gesture, music, and audio modalities so that the cross-modal claims of Chapter~\ref{ch:essay-crossmodal} can transition from illustrative examples to measured results, running larger-scale versions of the Recoverability experiment to establish statistical confidence intervals for collision rates under different evolutionary conditions, and extending the Reconstruct experiment to test diagnostic predictions from Corollary~4.1 by deliberately constructing histories with known reachable-set relationships and measuring whether reconstruction uncertainty aligns with the theorem's predictions. The second category includes category-theoretic reformulation of the admissibility relation to make functorial properties explicit, investigation of hierarchical control structures in which higher-level reference conditions influence lower-level admissibility constraints as discussed in the perceptual-control literature, and extension of the framework to multi-scale systems where admissibility relations at one temporal or organizational scale constrain but do not fully determine admissibility at other scales. Planned experiments in the repository that would test the framework under progressively more challenging conditions include the False Cognate Laboratory, which measures accidental resemblance rates and tests evidential requirements against false-positive thresholds; the Dialect Continuum, which challenges tree-based genealogical assumptions by modeling diffusion through continuous space; Borrowing Without Ancestry, which introduces horizontal transmission to test when contact is mistaken for common descent; the Babel Experiment, which investigates whether linguistic differentiation emerges from initially uniform populations without encoding differentiation as an explicit outcome; Maximum Language, which complements the implemented Minimum Language by beginning with excessive complexity and measuring what survives imperfect transmission; Glyph Evolution, which extends the framework to writing-system evolution; ASL Handshape Drift, which tests whether the admissibility apparatus remains useful when linguistic evolution occurs in a different articulatory and perceptual modality; The Last Similar Thing, which compares recency-based transmission against retrieval based on structural and contextual similarity; and Language Earth, the long-term integration experiment combining migration, contact, divergence, borrowing, innovation, writing, prestige, isolation, and reconnection while retaining complete ground truth for subsequent reconstruction attempts. These experimental extensions would stress-test the distinction between inference error and genuine information loss under conditions including sparse evidence, convergent evolution, modality change, and non-tree genealogy, addressing not merely whether reconstruction succeeds but why it fails when it does. The framework's implications for other domains involving inverse problems under lossy projection may prove relevant to fields as distinct as paleontology, archaeology, and signal processing, though such applications remain speculative absent domain-specific validation of the admissibility structure's appropriateness.

\chapter{Conclusion}
\label{ch:essay-conclusion}

The four theorems examined in this essay establish that representation, coherence, repair, and reconstruction are not four independent subjects requiring separate theoretical foundations but four aspects of a single admissibility relation examined at different points in its operation. Representation Invariance addresses the question of when physically distinct substrates count as semantically equivalent, answered by checking whether the admissibility structures they induce are related by an appropriate isomorphism. Admissibility Preservation addresses the question of what makes a trajectory coherent, answered by verifying that the reachable set remains nontrivial at each step rather than degenerating to emptiness or a predetermined singleton. Monotonic Relaxation addresses the question of when inadmissibility decreases, answered conditionally for trajectories following repair dynamics that increase reachable-set volume. Non-Identifiability addresses the question of when reconstruction becomes impossible in principle, answered by showing that observation maps factoring through reachable-set structure are generically non-injective whenever history space exceeds what that structure encodes.

What unifies these four results is that each derives from the same underlying apparatus: the admissibility relation $x \adm{R} y$ determining which continuations remain admissible from any given state, and the reachable set $\reach{x}$ collecting those admissible continuations into a structure that can be examined, compared, and projected. The formal development in \emph{Admissible Reconstruction} establishes that no additional primitive vocabulary is required beyond this minimal starting point. Representation equivalence is admissibility-structure isomorphism. Coherence is reachable-set nondegeneracy. Repair is reachable-set volume increase. Non-identifiability is equivalence-class structure under projection. The economy is not coincidental but reflects the fact that a single relation, when examined from different perspectives, generates all four phenomena as consequences.

The cross-modal investigation developed here applies this unified apparatus to the projection problem: determining what structural relations survive when latent intentional states are realized through different physical channels. Speech, gesture, music, and audio differ in which aspects of structure each modality preserves and which it necessarily discards, but the formal question remains the same across modalities. Does the projection preserve admissibility structure completely, as required by Theorem~1.1 for representational equivalence, or does it introduce distortions that break the isomorphism? Does the projected trajectory maintain nontrivial reachable sets at each step, satisfying the coherence condition of Theorem~2.1, or does it degenerate? If the system undergoes repair, does the trajectory follow the volume-increasing dynamics required by Theorem~3.1 for monotonic relaxation? And when an observer attempts to reconstruct latent structure from observable projection, does the observation map produce collisions as predicted by Theorem~4.2, and if so, are those collisions correctable or structural as diagnosed by Corollary~4.1?

The deepest result concerns the limits of reconstruction. Non-identifiability establishes that when distinct histories produce the same observable outcomes, the failure to distinguish them is not a weakness of the observer's inferential methods but a property of the observation channel itself. Two histories that share the same reachable-set structure have become, from the perspective of any observation map preserving only continuation structure, identical in the observable projection. The information required to distinguish them has not merely become difficult to access; it has ceased to exist in the projection. This is what it means for observables to be lossy: not that some details are omitted while others remain, as though careful observation could in principle recover the missing information, but that entire dimensions of variation in the history space collapse to equivalence classes in the observable space, and no refinement of observational technique operating within that space can recover the collapsed dimensions.

The diagnostic corollary provides the criterion for recognizing when such collapse has occurred. Given two candidate histories producing similar observables, check whether their reachable sets coincide. If they differ, the reconstruction uncertainty is correctable, and additional evidence could in principle resolve it. If they are equal, the uncertainty is structural, and no observation map factoring through reachable-set structure can eliminate it. This transforms reconstruction failure from an undifferentiated acknowledgment that evidence may be insufficient into a structured diagnosis with decidable answers: the failure is either correctable or fundamental, and the distinction can be determined by examining the reachable-set structure directly.

The framework developed here therefore addresses the projection problem not by asking whether reconstruction succeeds or fails in particular cases but by establishing what structural properties determine success and failure. Representations are equivalent when admissibility-isomorphisms exist. Trajectories are coherent when reachable sets remain nontrivial. Repair succeeds when reachable-set volume increases. Reconstruction is impossible when distinct histories share reachable-set structure. These are not four separate claims requiring independent justification but four views of one relation, and recognizing them as such clarifies both what the admissibility apparatus establishes and what it leaves for empirical investigation to determine.


\backmatter

\begin{thebibliography}{9}

\bibitem{lamport1978}
Leslie Lamport, ``Time, Clocks, and the Ordering of Events in a
Distributed System,'' \textit{Communications of the ACM}, 1978.

\bibitem{chandy1985}
K.~Mani Chandy and Leslie Lamport, ``Distributed Snapshots: Determining
Global States of Distributed Systems,'' \textit{ACM Transactions on
Computer Systems}, 1985.

\bibitem{shannon1948}
Claude E.~Shannon, ``A Mathematical Theory of Communication,'' \textit{Bell
System Technical Journal}, 1948.

\bibitem{cover2006}
Thomas M.~Cover and Joy A.~Thomas, \textit{Elements of Information
Theory}, 2nd ed., Wiley, 2006.

\bibitem{friston2010}
Karl Friston, ``The Free-Energy Principle: A Unified Brain Theory?,''
\textit{Nature Reviews Neuroscience}, 2010.

\bibitem{aubin1991}
Jean-Pierre Aubin, \textit{Viability Theory}, Birkh\"auser, 1991.

\bibitem{powers1973}
William T.~Powers, \textit{Behavior: The Control of Perception}, Aldine,
1973.

\end{thebibliography}


\end{document}
